\chapter{Lifting in every positive characteristic}\label{tx:sec:odd}
\label{tx:ch:prime}
\section*{The property we need}
Recall the \emph{block operator} used to construct the lift:
\[
                       z=\sum_i c_i(\sigma_i-I).
\]
Here $\sigma_i$ translates the added block in its $i$th coordinate,
$I$ is the identity operator, and the $c_i$ are the chosen field
coefficients. Thus $z$ is a linear operator on our labelled vector
spaces, not a scalar. The equation $z^2=0$ says that applying $z$
\emph{twice} gives zero: $z(z(v))=0$ for every vector $v$. It does
not require $z(v)=0$ on the first application.

For a concrete example over $\mathbb F_2$, subtract the identity
from the operator which swaps two coordinates. The resulting matrix is
\[
 z=\begin{pmatrix}1&1\\1&1\end{pmatrix}\ne0,
 \qquad z^2=\begin{pmatrix}0&0\\0&0\end{pmatrix}.
\]
Each entry in its square is $1+1=0$ in the field. Such an operator is
called \emph{square-zero}; it sends its own image to zero.

The characteristic-two proof used two special properties of the
general block operator: it was square-zero, and its matrix on the
permutation basis equalled its transpose. The latter holds because
each coordinate translation is an involution, so its permutation
matrix is symmetric. In odd characteristic these properties require
different replacements. The combinatorial core, however, used signed
sums and triangular elimination without division by a group order,
and remains available.

This chapter proves every replacement needed for the same linear
retained-label budget. The block has $q=p^a$ labels and $z^p=0$.
The deletion quotient must discard every nilpotent chain longer than
one, not just the full length-$p$ chains. A common witness for the
top power $u=z^{p-1}$ yields a multi-factor projector for that actual
quotient. We first prove the requisite purity on the tuple generators,
then obtain identities in their common annihilator algebra, and only
then apply those identities to arbitrary subspaces and quotients.

The layer splitting now has $p$ levels. Duality uses the block
relabelling $v\mapsto-v$ to handle transpose. Finally, the same rook
realisation repairs the top-power residual without new labels.
These steps supply the exact-deletion theorem in every prime
characteristic.

The conclusion has the same progression form as before: fix $d$,
choose $q=p^a$ larger than its support budget, and lift with degree $d$
at every step. Odd characteristic changes the algebra which builds the
lift, not the meaning of repeated lifting or the linear label cost.

\section{The block operator in prime characteristic}
A linear operator is \emph{nilpotent} when some number of successive
applications gives zero. A nilpotent chain has independent members
$c,zc,\ldots,z^{j-1}c$, with $z^jc=0$; applying $z$ moves one step
down this chain. These are the elementary blocks often called Jordan
chains. The equation $z^p=0$ bounds their length by $p$, but does not
by itself exclude lengths between $2$ and $p-1$. That exclusion is
a property to prove for the actual block action.

Let now $\operatorname{char}K=p$ for any prime. Keep $N=n+q$ for the larger universe, with $n$ retained labels and $q$ added block labels. The support budget $r$ bounds the tuple lengths needed by the proof diagram; it is not the block exponent $a$ or the proof degree $d$.
As in the characteristic-two construction, we temporarily extend the
coefficient field if needed and write $K$ for the working field below.
The descent in Lemma~\ref{tx:lem:field} preserves the original degree
question. Take a regular block $A=\Fp^a$ of size $q=p^a>r$, choose
$c_1,\ldots,c_a$ independent over $\Fp$ in $K$, and put
\[z=\sum_i c_i(\sigma_i-1).
\]
The generators commute and have order $p$, so the binomial formula in characteristic $p$ gives
\[z^p=\sum_i c_i^p(\sigma_i^p-1)=0.
\]
For $p>2$ we cannot replace this by $z^2=0$. For example, for a three-cycle $\sigma$ in characteristic three,
\[z=\sigma-I,\qquad z^2=\sigma^2+\sigma+I\ne0
\]
on its three-point permutation space.

The binomial step uses the elementary characteristic-$p$ identity
$(U+V)^p=U^p+V^p$ for \emph{commuting} operators: each intermediate
binomial coefficient is divisible by $p$. It is valid for the coordinate
translations because they commute. It is not a rule for arbitrary
matrix sums, nor a statement that powers of unrelated operators commute.

\subsection{The deletion quotient}
Since $\im z$ need not be contained in $\ker z$, the characteristic-two expression $\ker z/\im z$ is no longer defined. Use
\begin{equation}\label{tx:eq:Dp}
\D(M)=\ker z/(\ker z\cap\im z).
\end{equation}
A nilpotent chain has basis $v_0,\ldots,v_{j-1}$ with $zv_0=0$ and $zv_i=v_{i-1}$ for $i\ge1$. Its kernel is $Kv_0$. If $j=1$, its image is zero and the chain survives in \eqref{tx:eq:Dp}. If $j\ge2$, $v_0$ is in the image and the chain disappears. Thus deletion keeps the one-dimensional chains and discards every longer chain.

One might instead try $\ker z/\im z^{p-1}$. It is a defined quotient, but it is the wrong one: a chain of length two in characteristic three would survive because $z^2=0$ there, even though it should be discarded by \eqref{tx:eq:Dp}. Keeping the quotient precise is important.

\section{The commutator argument in characteristic \texorpdfstring{$p$}{p}}\label{tx:sec:commutatorp}
Work on one moved orbit of a permutation basis. Choose a basis
object $x_0$, and let $H=\{v\in A:vx_0=x_0\}$ be its stabiliser.
Because the orbit is moved, $H$ is a proper subspace of $A$.
Choose a nonzero linear functional $\ell:A\to\Fp$ vanishing on
$H$, and define $\pi(vx_0)=\ell(v)$. This is well defined because
two translations taking $x_0$ to the same object differ by an
element of $H$. Let $D$ multiply each basis vector by its orbit
coordinate $\pi$, and let $S_i$ act by translation in coordinate
$i$. With $\varepsilon_i$ the $i$th unit vector, put
$e_i=\ell(\varepsilon_i)$ and $\alpha=\sum_i c_ie_i\ne0$.
The choices are made orbit by orbit. Exactly as before,
\[Dz-zD=W=\sum_i c_ie_iS_i,\quad Wz=zW,\quad W^p=\alpha^pI.
\]
The commutator $Dz-zD$ compares the order of two operations.
Here $W$ names that commutator operator, distinct from a polynomial
ambient space $W_N$ or a consequence space $W_d(n)$.
Its value $W$ is invertible for a concrete reason:
$W^{-1}=\alpha^{-p}W^{p-1}$. Independence of the $c_i$ makes the
scalar $\alpha$ nonzero. Thus it is $W$ whose inverse is used here,
not the nilpotent operator $z$.
Consequently the auxiliary linear operator $B=W^{-1}D$ satisfies
\begin{equation}\label{tx:eq:commutator}[B,z]=Bz-zB=I.\end{equation}
Here $B$ names an operator, not the set of block labels used in the
earlier chapters; brackets denote the displayed commutator of
composed operators.
This proof works also on a moved permutation basis tensored with an arbitrary representation: $D$ reads only the basis coordinate, while the $S_i$ act diagonally and still have order $p$.

There is an elementary model for this commutator. On the space with
basis $1,T,\ldots,T^{p-1}$, work modulo $T^p$, let $z_0$ multiply
by $T$, and let $B_0$ differentiate. Differentiation is well defined
modulo $T^p$ because $(T^p)'=pT^{p-1}=0$. For $j<p-1$,
$(B_0z_0-z_0B_0)T^j=((j+1)-j)T^j=T^j$.
At $j=p-1$ the first term is zero and the second has coefficient
$-(p-1)=1$ in the field. Hence $[B_0,z_0]=I$ on the whole space.
The actual orbit construction above produces its commutator directly;
this model explains why a finite-dimensional commutator can equal the
identity in characteristic $p$.

If $zm=0$, commuting $z$ past powers of $B$ yields
\[zB^jm=-jB^{j-1}m,
\qquad z^jB^jm=(-1)^jj!\,m.
\]
The first identity follows by induction from $zB=Bz-I$; applying it repeatedly gives the second. At $j=p-1$ the scalar $(p-1)!$ is nonzero in characteristic $p$, because every factor $1,\ldots,p-1$ is nonzero. Therefore
\begin{equation}\label{tx:eq:freep}\ker z=\im z^{p-1}
\end{equation}
on the moved summand. The reverse inclusion uses $z^p=0$.

The displayed formula starts with a kernel vector $m$ and constructs
a preimage at the full height $p-1$ by applying $B^{p-1}$ and dividing
by $(-1)^{p-1}(p-1)!$. Thus every vector killed by $z$ on a moved
orbit already comes from the top power. A full chain is a copy of
the truncated polynomial model just described; a direct sum of such
chains is what the earlier phrase ``free for $z$'' means here.
This is a statement about the block operator, not freeness for the
full symmetric-group action.
Equivalently, let the variable $T$ in the ring $K[T]/(T^p)$ act
as $z$. A full chain generated by $c$ has basis
$c,zc,\ldots,z^{p-1}c$ and is one free copy of that ring.
A direct sum of full chains is therefore a free module for this
truncated polynomial ring. This is the precise module meaning of
freeness used in the calculation.

This is the full replacement for the square-zero freeness calculation. It does not divide by $p$, by $|A|$, or by a symmetric-group order. It divides only by a nonzero orbit coefficient and by $(p-1)!$.

On a bounded-support permutation basis the retained summand has $z=0$ and the moved summand satisfies \eqref{tx:eq:freep}. Therefore the identifications with the retained basis and the tensor identifications from Chapter~\ref{tx:sec:block-algebra} still hold, with quotient \eqref{tx:eq:Dp}. Stable matrices are transported by exactly the same retained-component argument.

\subsection{A three-cycle computation}
In characteristic three, let $\sigma$ cycle $e_0,e_1,e_2$.
The operator $z=\sigma-I$ has
\[
 z^2=I+\sigma+\sigma^2,\qquad
 z^2 e_0=e_0+e_1+e_2,\qquad z^3=0.
\]
Its kernel is the line spanned by $e_0+e_1+e_2$, and this line is
exactly $\im z^2$. It lies in $\im z$, so the deletion quotient
is zero on this moved orbit. A separate fixed basis vector has
zero image and survives. This is the prime-characteristic counterpart
of the moved two-point calculation, with a full chain of three
levels instead of two.

The commutator argument proves the same conclusion for every moved
orbit, including those with nontrivial stabilisers and those tensored
with an arbitrary representation. The factorial is $(p-1)!$, not
$p!$: each factor being divided by is nonzero. It is this explicit
calculation that excludes shorter nontrivial nilpotent chains in
the original moved tuple modules.

\section{The top-power witness}\label{tx:sec:topwitness}
Set $m=p-1$ and $u=z^m$. Here $m$ records this fixed power
exponent, rather than a tuple length or support size. Like $z$, $u$ is a linear operator; $u^2=0$ because
$2(p-1)\ge p$ and $z^p=0$. Recall that $J_r(N)$ consists of the
group-algebra elements acting as zero on all ordered-tuple spaces
$P_b(N)$ with $b\le r$. The common certificate becomes
\begin{equation}\label{tx:eq:topwitness}uwu-u\in J_r(N),\qquad\text{or }uwu=u\text{ on every }P_b(N),\ b\le r.
\end{equation}
This says that $w$ takes the bottom of a full length-$p$ chain back to a preimage at its top. It is enough to build a projector for the actual quotient, but more steps are needed than in characteristic two.

\subsection{Fixed directions, full chains, and purity}
The condition called \emph{purity} below is
$\ker z\cap\im z\subseteq\im z^{p-1}$. It says that a kernel
vector which already has a preimage under $z$ also has a preimage
under its top power. On a chain of length $j\ge2$, the bottom vector
is in $\ker z\cap\im z$. It is in $\im z^{p-1}$ only when the
chain has the full length $p$. On a fixed line the intersection is
zero. This is why purity distinguishes fixed directions and full
chains from shorter nontrivial chains.

\begin{definition}[Fixed and full-chain decomposition]\label{tx:def:full-chains}
For an operator $z$ with $z^p=0$, a fixed and full-chain decomposition is
\[
             V=F\oplus\bigoplus_{j=0}^{p-1}z^jC,
             \qquad zF=0,
\]
where the sum is direct and $z^{p-1}:C\to z^{p-1}C$ is injective.
Thus each basis vector $c$ of $C$ supplies the full chain
$z^{p-1}c,\ldots,zc,c$; $F$ consists of fixed directions.
\end{definition}
``Fixed'' here means killed by $z$. It need not mean fixed by every
label permutation. The direct-sum decomposition is a statement about
this nilpotent operator, not an equivariant decomposition for the
full symmetric group.

This definition will describe the layer splitting constructed below.
For the original tuple modules, purity follows directly from the orbit
calculation, without choosing chain bases:
\begin{equation}\label{tx:eq:purity}\ker z\cap\im z\subseteq\im z^{p-1}.
\end{equation}
The image has no retained component, and a moved kernel vector is in
the top-power image by \eqref{tx:eq:freep}. The algebraic contract is
therefore: prove purity on these generators; use one top-power witness
to obtain projector identities there; transfer the identities through
$J_r(N)$ to prove exactness on the whole annihilator class. A chosen
chain complement need not be compatible with submodules or quotients.

\begin{example}[Why the distinction matters]
In characteristic three, a lone length-two chain has $u=z^2=0$. The relation $uwu=u$ is then vacuous for any $w$, but $\ker z\cap\im z$ is nonzero and $\im u=0$. Such a chain does not satisfy \eqref{tx:eq:purity}. The argument must establish its projector identities on the actual tuple modules first; the top-power relation alone on an arbitrary vector space is insufficient.
\end{example}

\subsection{The ordered projector}
For $p>2$, the square-zero operator $u=z^{p-1}$ can have a larger
kernel than $z$. In characteristic three, on the chain $zv_2=v_1$,
$zv_1=v_0$, $zv_0=0$, the top power kills $v_1$, while $z$ does not. If $w$
sends $v_0$ to $v_2$ and kills the other two vectors, the
characteristic-two expression $(I-uw)(I-wu)$ would leave $v_1$
unchanged. It therefore cannot model the deletion quotient for $z$.
The ordered product below uses intermediate powers of $z$ to remove
every level, not merely the top and bottom ones.

For a vector satisfying $z^{h+1}v=0$, call $h$ an upper bound on
its chain height. This height is independent of polynomial degree.
The proof controls the height of the image after each successive
factor. The permitted height decreases; the degree budget of a PC
proof is not being decreased. Only the completed product is claimed
to be the required projector; its factors are not assumed to commute.

\begin{lemma}[Projector from purity and a top-power witness]\label{tx:lem:prime-projector}
Let $z,w$ be endomorphisms of a vector space over a field of
characteristic $p$, and put $m=p-1$. Suppose $z^{m+1}=0$,
$z^mwz^m=z^m$, and \eqref{tx:eq:purity} holds. Define
\begin{equation}\label{tx:eq:prime-projector}
S_0=I,\qquad S_{k+1}=(I-z^kwz^{m-k})S_k\quad(0\le k<m),
\qquad e=(I-z^mw)S_m.
\end{equation}
Then $ze=ez=0$, $e^2=e$, and $ev=v-z^mwv$ whenever $zv=0$.
\end{lemma}
\begin{proof}
Operator products act from right to left. Their order is part of the
construction; no factor is commuted past $w$.

\paragraph{Eliminating high levels.}
At step $k$ put $h=m-k$. Inductively $z^{h+1}S_k=0$; this is
true at $k=0$ because $z^{m+1}=0$. For $k<m$ we have $h\ge1$, so
$y=z^hS_kv$ belongs to $\ker z\cap\im z$, hence to $\im z^m$.
The witness gives $z^mwy=y$, and therefore
\[
 z^hS_{k+1}v
       =z^h(I-z^{m-h}wz^h)S_kv=y-z^mwy=0.
\]
The permitted height drops from $h$ to $h-1$ at each step. Thus
$z^{m+1-k}S_k=0$ for every $k$, and finally $zS_m=0$.
As $z^{m+1}=0$, the last factor also gives $ze=0$.

\paragraph{Action on a kernel vector.}
If $zv=0$, every factor in $S_m$ fixes $v$, because its rightmost power of $z$ has exponent at least one. Therefore
\begin{equation}\label{tx:eq:ekernel}ev=v-z^mwv\quad(zv=0).
\end{equation}

\paragraph{Killing the image.}
The first factor $I-wz^m$ fixes $\im z$, since $z^{m+1}=0$. All later factors preserve $\im z$ because their correction terms have a left factor $z^k$ with $k\ge1$. Hence $S_mzv\in\ker z\cap\im z\subseteq\im z^m$. The final factor $I-z^mw$ kills this vector by the witness relation. Thus $ez=0$.

\paragraph{Idempotence.}
The witness gives
\[z^mwe=(z^mw-z^mwz^mw)S_m=0.
\]
Since $ev\in\ker z$, formula \eqref{tx:eq:ekernel} now gives $e^2v=ev$. We have proved $ze=ez=0$ and $e^2=e$.
\end{proof}

The final factor kills the bottom of each full chain as well.
Formula \eqref{tx:eq:ekernel} is especially useful for lifting:
only the final correction acts on a kernel vector. It is the identity needed for
$\Lift=e\iota$ and the retained-coordinate right inverse.

\subsection{Passing from tuples to every relevant space}
Recall that $\mathcal C_r(N)$ contains every finite-dimensional
label representation killed by $J_r(N)$, including the invariant
subspaces and quotients used by the consequence calculation. The
point of a common witness is that identities verified on all the
tuple generators then hold on every space in this class.

Concretely, $K[S_N]/J_r(N)$ is the algebra in which two permutation
combinations are identified when their difference kills every tuple
test space of length at most $r$. The action on each $M\in\mathcal C_r(N)$ factors
through this quotient. Consequently an identity such as $e^2=e$
proved on all the tuple spaces is an identity of this algebra and
acts on every such $M$. A chosen chain complement need not restrict
to a submodule, but an identity of common operators does. The argument
therefore transfers identities rather than separately chosen bases.

Naturality here concerns maps $f$ of label representations: they
intertwine every common group-algebra operator, giving
$f(ev)=e f(v)$. This does not assert that $e$ commutes with every
label permutation. Its image is used as a vector-space model for
the deletion quotient, which carries the appropriate residual action.

For the actual block action, purity on every tuple generator has
already been proved in \eqref{tx:eq:purity}. The proposition below
isolates the implication from that property and a common witness to
exactness. Its purity premise records the input to this algebraic
step; it is not an outstanding assumption of the block construction.

\begin{proposition}[Common identities give exact deletion]\label{tx:prop:prime-exactness}
For the block operator $z$ above, $z^p=0$. Suppose purity holds on
every $P_b(N)$, $b\le r$, and one
$w\in K[S_N]$ satisfies \eqref{tx:eq:topwitness}. The group-algebra
expression $e$ in \eqref{tx:eq:prime-projector} gives a natural
vector-space isomorphism $\D(M)\simeq eM$ for every
$M\in\mathcal C_r(N)$, and $\D$ is exact on this class.
\end{proposition}
\begin{proof}
The lemma proves the identities on the tuple generators, so $ze$, $ez$
and $e^2-e$ belong to $J_r(N)$. They act as zero on every object of
$\mathcal C_r(N)$, whether or not that object has a fixed and full-chain
decomposition. Formula \eqref{tx:eq:ekernel} follows there directly from
the ordered product, since all factors in $S_m$ fix a kernel vector.

The quotient-to-image map sends $[v]$ to $ev$. It is well defined
because $ez=0$, and $eM\subseteq\ker z$ because $ze=0$.
For $zv=0$, the difference $v-ev=z^{p-1}wv$ belongs to both
$\im z$ and $\ker z$. Taking the quotient class of $ev$ therefore
returns $[v]$. Conversely, idempotence makes every element of $eM$
fixed by $e$. These maps are inverse for the actual quotient
\eqref{tx:eq:Dp}.

Every module map commutes with the same group-algebra expression $e$.
For an exact sequence $0\to U\xrightarrow{i}V\xrightarrow{f}W\to0$,
restriction to $eU\to eV$ preserves injectivity. If $y\in eW$, any
lift $v\in V$ gives the projected lift $ev$, since $f(ev)=ef(v)=y$.
If $v\in eV$ maps to zero, write $v=i(u)$ and then
$v=ev=i(eu)$. Thus the projected sequence is exact. The natural
quotient-to-image isomorphisms prove exactness of deletion, without a
separate purity assumption on subobjects. As in
Chapter~\ref{tx:ch:square-zero}, $eM$ is a vector-space model;
the quotient retains the residual label action.
\end{proof}

\begin{example}[A full chain in characteristic three]
Let $zv_2=v_1$, $zv_1=v_0$, $zv_0=0$, and let $wv_0=v_2$, with $w$ zero on the other chain vectors. Then $z^2wz^2=z^2$. The projector is
\[e=(I-z^2w)(I-zwz)(I-wz^2).
\]
The rightmost factor removes $v_2$, the middle factor removes $v_1$, and the leftmost removes $v_0$. If a separate fixed vector $v_{\rm fix}$ is added and $wv_{\rm fix}=0$, every factor fixes $v_{\rm fix}$. Thus the projector removes exactly the full chain and retains the fixed direction. For $p=2$ the same formula has only the two factors $(I-zw)(I-wz)$ from Chapter~\ref{tx:sec:block-algebra}.
\end{example}

\section{The linear construction in odd characteristic}\label{tx:sec:oddlinear}
We now adapt the characteristic-two construction from
Chapter~\ref{tx:ch:linear-layers}, using characteristic-independent
combinatorial ingredients and new operator arguments. The combinatorial
ingredients are the partition of unity in Section~\ref{tx:sec:partition}, the cone
and signed component lifts in Section~\ref{tx:sec:cone}, and the
rook forcing and permutation realisation in
Chapter~\ref{tx:ch:rooks}. Their insertion and cancellation formulas
use $1$ and $-1$ without a characteristic-two cancellation, so these
combinatorial proofs remain valid in every field. They supply
the same retained-label count $2b-1$.

The operator arguments do require replacements. The two-level
splitting must become a full $p$-level splitting, and the
characteristic-two identity $z^{\mathsf T}=z$ must be replaced by
conjugacy with the transpose. The following subsections establish
these replacements before using the common correction contract
with $u=z^{p-1}$.

Fix the current tuple length $b\le r$, and recall
$V=P_b(N)$, $V'=\sum_i\im\partial_i^{\mathsf T}$,
$\mathcal Z_b=\bigcap_i\ker\partial_i$ and $T_b=V/V'$.
Thus $V'$ is the incoming shorter-tuple space, $\mathcal Z_b$ is
the zero-marginal space, and $T_b$ is the quotient by the incoming
directions. The filtration $\mathcal Z^{\le\ell}$ consists of
zero-marginal vectors supported on tuples with at most $\ell$ block
positions. Its prospective quotient layer
is
\[
G_\ell=\bigoplus_{|J|=\ell}
       K[\Inj(J,A)]\otimes\mathcal Z_{[b]\setminus J}(n),
\]
where $J\subseteq[b]$ is a set of positions and $A$ is the block-label set.
These are the same layer spaces used in the characteristic-two
construction; only the block action and its chain structure change.

\subsection{Constructing a fixed and full-chain decomposition}
A nonempty block pattern has an image set of size $\ell<q$.
Give it a grade $e\in\Fp$ using the orbit coordinate of that image
set, as in Chapter~\ref{tx:ch:linear-layers}. Let $E$ be the space on
all these block patterns, and let $E_e$ span the
patterns of grade $e$, and let $D$ read the grade. The same commutator
$W=[D,z]$ is invertible and commutes with $z$.
In this subsection $e$ is a scalar grade; the projector denoted by
$e$ in the preceding sections is an operator and plays no role in
the grade decomposition.

Fix a grade $e$ and suppose $\sum_{j=0}^{p-1}z^jv_j=0$ with
$v_j\in E_e$. The commutator identity gives
\[
 (D-eI)z^jv_j=jWz^{j-1}v_j\quad(j\ge1).
\]
Applying $D-eI$ to the relation and cancelling $W$ leaves
$\sum_{j=1}^{p-1}jz^{j-1}v_j=0$. Repeat $p-1$ times to obtain
$(p-1)!v_{p-1}=0$, hence $v_{p-1}=0$. Remove that term and repeat
with one fewer power. All $v_j$ vanish. This proves both directness
of $\sum_jz^jE_e$ and injectivity of each $z^j$ on $E_e$.
It does not require $W$ to preserve the grade: after each cancellation
the coefficients are still the original grade-$e$ vectors.

Consequently $p\dim E_e\le\dim E$ for each of the $p$ grades.
As they partition the basis, $\sum_e\dim E_e=\dim E$. Summing the
inequalities forces equality in every one. In particular
\[
                   E=\bigoplus_{j=0}^{p-1}z^jE_0.
\]
The grade depends only on the image set, so it is stable under
position permutations and splits by position domain. It ignores the
retained tuple. Tensoring each domain component with its retained
zero-marginal factor therefore gives
$G_\ell=\bigoplus_{j=0}^{p-1}z^jQ_\ell$, with $Q_\ell$
position-stable and $z^{p-1}$ injective on $Q_\ell$.

This is the $p$-level version of the two-level section construction.
The grade $e\in\mathbb F_p$ is a coordinate on an image-set orbit;
$\ell$ instead counts block positions. To make a section commute
with $z$, choose preimages of the chain starts $q\in Q_\ell$,
then prescribe the preimage of $z^jq$ to be $z^j\sigma(q)$.
Directness and injectivity make all the chain coordinates unique.
The last step goes to zero because $z^p=0$, so the prescription
is compatible with $z$ at every level without an averaging operation.

For a layer quotient map $f$ and the position-equivariant section
$\sigma$ supplied by the signed lift, define
\[
 \tau\Bigl(\sum_{j=0}^{p-1}z^jq_j\Bigr)
                  =\sum_{j=0}^{p-1}z^j\sigma(q_j).
\]
Uniqueness of the decomposition makes the formula well defined;
$f\tau=I$ follows from $fz=zf$. Commuting with $z$ is checked by
shifting each power up by one; the term at power $p$ is zero.
Its image splits that filtration quotient as an actual subspace.
Position stability of $Q_\ell$ and position equivariance of $z$ and
$\sigma$ also make $\tau$ position-equivariant.
Successively splitting the quotients gives
\begin{equation}\label{tx:eq:socsplitp}
\mathcal Z_b(n+q)=\mathcal Z_b(n)\oplus\bigoplus_{j=0}^{p-1}z^jC,
\end{equation}
where $C$ is position-stable and $z^{p-1}$ is injective on $C$.
This is a fixed and full-chain decomposition in the sense of
Definition~\ref{tx:def:full-chains}. The hypotheses remain
$n\ge2b-1$ and $b<q$: making the sections commute with $z$ adds no
labels to the retained-label cost of the signed lift.

\subsection{Top and zero-marginal spaces}
Write $V=P_b(N)$. We verify the four properties used in
Lemma~\ref{tx:lem:layer-interface}, now with $u=z^{p-1}$ in place of
$z$. In \eqref{tx:eq:socsplitp}, every vector has a unique expression
$s+\sum_{j=0}^{p-1}z^jc_j$, with $s\in\mathcal Z_b(n)$ and
$c_j\in C$. Since $zs=0$,
\[
 z\Bigl(s+\sum_{j=0}^{p-1}z^jc_j\Bigr)
                         =\sum_{j=0}^{p-2}z^{j+1}c_j.
\]
Directness and injectivity of the powers of $z$ on $C$ force
$c_0=\cdots=c_{p-2}=0$ in a kernel vector. Hence
\[
 \ker(z|_{\mathcal Z_b})=\mathcal Z_b(n)\oplus uC,
                         \qquad u\mathcal Z_b=uC.
\]
If a vector belongs to $uV\cap\mathcal Z_b$, it is therefore $s+uc$.
Both that vector and $uc$ have zero retained coordinates: $u$ kills
retained tuples and preserves the moved-tuple summand. Thus $s=0$, and
\[
                         uV\cap\mathcal Z_b=u\mathcal Z_b.
\]
Define $w_C$ on $\mathcal Z_b$ by sending $uc$ to $c$ and killing
the retained summand and the levels $C,zC,\ldots,z^{p-2}C$.
Injectivity of $u$ on $C$ makes this well defined; position stability
makes it position-equivariant. It satisfies $uw_Cu=u$ there.

\subsection{Duality with the actual transpose operator}
The transpose changes which space carries the inverse used by
duality. In the tuple pairing, $\langle uv,w\rangle
=\langle v,u^{\mathsf T}w\rangle$. Thus the operator induced by
$u$ on the quotient $T_b$ is paired with $u^{\mathsf T}$ on
$\mathcal Z_b$, rather than automatically with $u$ itself.
This is the part of the characteristic-two argument that must change.

All transposes here are ordinary bilinear matrix transposes. They do
not conjugate the coefficients $c_i$. A permutation matrix transposes
to the matrix of its inverse, giving
\[
 z^{\mathsf T}=\sum_i c_i(\sigma_i^{-1}-1).
\]
The map
$\iota(v)=-v$ on the block is a genuine label permutation, fixes
retained labels and commutes with position permutations. It therefore
conjugates the entire tuple-layer construction for $z$ to that for
$z^{\mathsf T}$: $\iota\sigma_i\iota^{-1}=\sigma_i^{-1}$ and
$\iota u\iota^{-1}=u^{\mathsf T}$. It preserves $V'$ and
$\mathcal Z_b$, since these are defined by label-equivariant marginals.
The symbol $\iota$ here denotes this block inversion and its action
on tuple spaces; it is different from the inclusion
$\iota_{\mathcal Z}:\mathcal Z_b\hookrightarrow V$ used in the
correction formulas.

For a three-element block labelled by $\mathbb F_3$, inversion fixes
$0$ and swaps $1,2$. Conjugating translation by $1$ with this swap
gives translation by $-1$. This is the elementary meaning of
$\iota z\iota^{-1}=z^{\mathsf T}$: the change is an actual
relabelling of the block, so the marginal constructions remain valid.
The inverse-image notation $(u^{\mathsf T})^{-1}(U)$ below means
$\{v:u^{\mathsf T}v\in U\}$, not an inverse matrix for $u^{\mathsf T}$.

In particular,
$u^{\mathsf T}V\cap\mathcal Z_b\subseteq
u^{\mathsf T}\mathcal Z_b$. This gives
\[
 (u^{\mathsf T})^{-1}(\mathcal Z_b)
                   =\ker u^{\mathsf T}+\mathcal Z_b.
\]
Taking orthogonal complements in the nondegenerate tuple pairing,
\[
 (uV')^\perp=(u^{\mathsf T})^{-1}(\mathcal Z_b),\qquad
 (uV\cap V')^\perp=\ker u^{\mathsf T}+\mathcal Z_b.
\]
Their equality proves $uV\cap V'=uV'$.

The perfect pairing $T_b\times\mathcal Z_b\to K$ identifies a
quotient vector with a linear functional on $\mathcal Z_b$.
The dual of a map $h$ sends such a functional $\lambda$ to
$\lambda\circ h$. This reverses composition order, as with ordinary
matrix transposes. In the following construction, this dual map is
$w_T$; the displayed pairing calculation verifies its inverse
condition on the correct space.

To specify the quotient witness, put
\[
 h=\iota w_C\iota^{-1}:\mathcal Z_b\longrightarrow\mathcal Z_b,
                  \qquad u^{\mathsf T}hu^{\mathsf T}=u^{\mathsf T}.
\]
The tuple dot product induces the perfect pairing
$\langle\ ,\ \rangle:T_b\times\mathcal Z_b\to K$.
Define $w_T:T_b\to T_b$ by
$\langle w_Tt,\zeta\rangle=\langle t,h\zeta\rangle$.
Then
\[
 \langle uw_Tut,\zeta\rangle
 =\langle t,u^{\mathsf T}hu^{\mathsf T}\zeta\rangle
 =\langle ut,\zeta\rangle,
\]
so $uw_Tu=u$. The pairing is invariant under position permutations,
and $h$ commutes with them; equivalently, their matrices transpose to
their inverses. Thus $w_T$ is position-equivariant. This supplies all
four correction properties with the actual transpose operator.

\subsection{A trace from transpose to correction in characteristic three}
The following small linear-algebra model displays the maps in the
interface. Let $V=Kf\oplus Ke_0\oplus Ke_1\oplus Ke_2$ over
$K=\mathbb F_3$, let $\sigma$ fix $f$ and cycle the three $e_i$, and
put $z=\sigma-I$, $u=z^2$. Choose the incoming space $V'=Kf$.
For the ordinary dot product its orthogonal complement is
$\mathcal Z=\Span\{e_0,e_1,e_2\}$, and $T=V/V'$ is paired perfectly
with $\mathcal Z$. These are model spaces for the interface, rather
than an assertion about a particular tuple length.

Write $v_{\Sigma}=e_0+e_1+e_2$. On $\mathcal Z$ the top power satisfies
$ue_i=v_{\Sigma}$, so its matrix is the all-one $3\times3$ matrix. The
chain-start space can be chosen as $C=Ke_0$: the vectors
$e_0,ze_0,v_{\Sigma}$ are independent. Define $w_C$ by
\[
 w_Ce_2=e_0,\qquad w_Ce_0=w_Ce_1=0.
\]
It sends $v_{\Sigma}$ back to $e_0$ and kills $C$ and $zC$, as the construction
requires. Block inversion fixes $e_0$ and swaps $e_1,e_2$; hence
$h=\iota w_C\iota^{-1}$ sends $e_1$ to $e_0$ and kills the other
two basis vectors. Taking the dual in the displayed pairing gives
\[
 w_T[e_0]=[e_1],\qquad w_T[e_1]=w_T[e_2]=0.
\]
For example, $uw_Tu[e_2]=uw_T[v_{\Sigma}]=u[e_1]=[v_{\Sigma}]=u[e_2]$.
The same calculation applies to every basis vector. This shows the
transpose reversal explicitly: $w_T$ is the dual of $h$, with its
specified domain $T$, rather than a reuse of $w_C$ on quotient vectors.

Now start this model's repair with $w_0=0$, so $\varepsilon=-u$.
The typed correction
$\varphi=\iota_{\mathcal Z}(-w_C\bar\varepsilon w_T)\pi$
sends $e_0$ to $e_0$ and kills $f,e_1,e_2$: indeed $w_T[e_0]=[e_1]$,
$\bar\varepsilon[e_1]=-v_{\Sigma}$, and $-w_C(-v_{\Sigma})=e_0$.
Therefore $u\varphi ue_i=v_{\Sigma}=-\varepsilon e_i$, while both sides
vanish on $f$. Thus the residual is cancelled on every vector.

Using $w=\varphi$ in the ordered projector makes the three stages
equally visible. The factor $I-wu$ kills $e_0$ and fixes $ze_0$ and
$v_{\Sigma}$; $I-zwz$ then kills $ze_0$ and fixes $v_{\Sigma}$; $I-uw$ kills $v_{\Sigma}$.
Every factor fixes $f$. The full chain is removed and the fixed line
survives. On actual tuple spaces the preceding construction supplies
position equivariance and the incoming/zero-marginal conditions;
the rook theorem supplies the common permutation combination.

\subsection{The same correction, with the correct power}
The remaining correction is the algebraic contract from
Chapter~\ref{tx:ch:linear-layers}, with $u$ in place of $z$.
Its typed composition uses the quotient $T_b$ at the input and the
zero-marginal space at the output. It uses neither an arbitrary choice
of quotient representatives nor a characteristic-two cancellation.
The preceding chain and transpose arguments have supplied its two
inverse conditions; the rook theorem then realises its matrix on the
same existing label set.

Suppose $w_0$ works on shorter tuples for the top-power relation. Its error is
\[\varepsilon=uw_0u-u.
\]
It maps into $\mathcal Z_b$ and kills $V'$, giving the factorisation
$\varepsilon=\iota_{\mathcal Z}\bar\varepsilon\pi$, where
$\bar\varepsilon:T_b\to\mathcal Z_b$. Write $u_T$ and
$u_{\mathcal Z}$ for the induced operators. We now verify both inverse
relations explicitly. The error's image lies in
$uV\cap\mathcal Z_b=u\mathcal Z_b$, so the witness on that
space gives $u_{\mathcal Z}w_C\bar\varepsilon=\bar\varepsilon$.
If $uv\in V'$, the
inclusion $uV\cap V'\subseteq uV'$ gives $uv=uv'$ for some
$v'\in V'$. As $\varepsilon=(uw_0-I)u$ and $\varepsilon V'=0$,
\[
                  \varepsilon v=(uw_0-I)uv'
                                  =\varepsilon v'=0.
\]
The quotient witness $u_Tw_Tu_T=u_T$ says that a representative
$v_T$ of $w_Tu_T\pi(v)$ satisfies $u(v_T-v)\in V'$. The preceding
implication therefore gives $\varepsilon v_T=\varepsilon v$.
Injectivity of $\iota_{\mathcal Z}$ gives
$\bar\varepsilon w_Tu_T\pi=\bar\varepsilon\pi$; surjectivity of
$\pi$ then gives $\bar\varepsilon w_Tu_T=\bar\varepsilon$.
No compatible choice of representatives is assumed.

With $\pi:V\to T_b$ and $\iota_{\mathcal Z}:\mathcal Z_b\to V$,
define the actual endomorphism
\[
       \varphi=\iota_{\mathcal Z}
                      (-w_C\bar\varepsilon w_T)\pi.
\]
Lemma~\ref{tx:lem:two-sided-correction}, applied with $u$ in place
of $z$, now gives $u\varphi u=-\varepsilon$. The map
is position-compatible, has image in $\mathcal Z_b$ and kills $V'$.
The characteristic-independent realisation theorem supplies
$j\in J_{b-1}(N)$ with this endomorphism. Hence
\[
        u(w_0+j)u-u=\varepsilon+uju=0
\]
on length $b$. On every shorter tuple module, $j=0$, so its old
witness relation is unchanged. The same label pool is used throughout.

Inducting on $b$ proves \eqref{tx:eq:topwitness} on the same label pool whenever $n\ge2r-1$. The longer projector gives exactness, the moved-orbit calculation gives natural basis and tensor deletion, and the original PC closure transport gives the same period statement.


\begin{proof}[Completion of the proof of Theorem~\ref{tx:thm:deletion}]
Begin at tuple length zero with $w_0=0$, since $z$ acts as zero on
constants. At length $b$, the preceding correction and
Theorem~\ref{tx:thm:realization} replace $w_{b-1}$ by $w_b$ satisfying
$z^{p-1}w_bz^{p-1}=z^{p-1}$ on every length at most $b$. The same
retained labels are used at every step. The largest requirement is
$\max_{1\le b\le r}(2b-1)=2r-1$; length zero needs no label.
The projector identities then imply exactness on $\mathcal C_r(N)$.
The orbit commutator calculation supplies the natural retained-basis
and tensor identifications, including the actual stable coefficient
maps. These arguments prove every assertion of Theorem~\ref{tx:thm:deletion}.
\end{proof}

\section*{Returning to the prescribed-target theorem}
The construction has proved the common top-power witness on every
tuple length up to an abstract budget $r$, at retained universe size
$n\ge\max\{0,2r-1\}$ and block size $q=p^a>r$. The tuple purity
calculation and the explicit projector prove exact deletion on the
full support class. The moved-orbit commutator proves natural basis
and permutation-first-factor tensor deletion. Stable coefficient
matrices therefore represent the actual smaller axiom and target
maps, and exactness transports the full PC closure.

For a fully invariant target the application uses $r=kd+s$. For
the indexed request it uses $r=2kd+s+t$, covering the certificate
source and the doubled polynomial pairing used to reflect target
membership. These choices are made by the target-transfer proof;
the finite algebra just proved is independent of the equations.
The target remains the prescribed full polynomial, including when
the degree budget is below its actual degree. No target is added
to the axioms and no semantic-entailment inference is introduced.

Finally, extend scalars only while constructing the generic block
operator and its common witness. Lemma~\ref{tx:lem:field} returns
actual derivability of each original-coefficient target to the original
field. Thus the construction applies to every coefficient field of
fixed positive characteristic. It has not asserted that a coefficientwise
projection of an arbitrary extension-field proof preserves each proof
line: the descent argument uses the finite-dimensional consequence
spaces and recovers a legal proof from their bases.

Both the onset and the period are linear in the fixed proof-degree
budget. Chapter~\ref{tx:chap:assemble-transfer} now assembles the actual
target-transfer theorem; the classification argument in
Chapter~\ref{tx:ch:profiles} inverts its two quantitative estimates,
with the exact easy-residue and finite-peak qualifications. The proof
engine is now complete in every prime characteristic.
